% Generated draft: requires mathematical and source review.
\documentclass[11pt]{article}
\usepackage[margin=25mm]{geometry}
\usepackage{fontspec}
\setmainfont{Latin Modern Roman}
\usepackage{amsmath,amssymb,mathtools,unicode-math}
\setmathfont{Latin Modern Math}
\usepackage{enumitem,needspace,xurl,hyperref}
\setlength{\parindent}{0pt}
\setlength{\parskip}{5pt}
\setlength{\emergencystretch}{3em}
\newcommand{\N}{\mathbb{N}}
\newcommand{\Z}{\mathbb{Z}}
\newcommand{\Q}{\mathbb{Q}}
\newcommand{\R}{\mathbb{R}}
\newcommand{\C}{\mathbb{C}}
\newcommand{\F}{\mathbb{F}}
\newcommand{\A}{\mathbb{A}}
\newcommand{\PP}{\mathbb{P}}
\newcommand{\E}{\mathbb{E}}
\newcommand{\eps}{\varepsilon}
\newcommand{\dd}{\,\mathrm{d}}
\begin{document}
\section*{UnsolvedMath ID: 30006419}
\textbf{Source problem code:} catalogue entry
\subsection{Definitions and mathematical statement}
Let \(S^2\) denote the standard round two-dimensional sphere. A metric space \((X,d)\) is quasiconvex if there exists \(L\geq 1\) such that every pair \(x,y\in X\) can be joined by a rectifiable curve of length at most \(L d(x,y)\).

The usual Sobolev filling formulation of a local quadratic isoperimetric inequality requires constants \(C>0\) and \(\ell_0>0\) such that every Lipschitz loop \(\gamma:S^1\to X\) with \(\operatorname{length}(\gamma)\leq\ell_0\) bounds a Sobolev disc \(v\in W^{1,2}(D,X)\) satisfying
\[
\operatorname{tr}(v)=\gamma,
\qquad
\operatorname{Area}(v)\leq C\operatorname{length}(\gamma)^2,
\]
where \(D\subset\R^2\) is the unit disc, \(W^{1,2}(D,X)\) denotes the metric-valued Sobolev class, and \(\operatorname{tr}\) denotes its boundary trace. This is a customary interpretation of the supplied hypothesis; the precise Sobolev and parametrized area conventions must be checked against the supplied references, since the record does not specify them.

The small-diameter hypothesis means that there exists \(\delta>0\) such that, for every continuous map \(f:S^2\to X\),
\[
\operatorname{diam}_d f(S^2)<\delta
\quad\Longrightarrow\quad
f\text{ is null-homotopic}.
\]
Here null-homotopic means that there exist \(x_0\in X\) and a continuous map \(H:S^2\times[0,1]\to X\) with \(H(z,0)=f(z)\) and \(H(z,1)=x_0\) for every \(z\in S^2\).

For a metric-valued Sobolev map \(u\), write \(\operatorname{ap\,md}u_z\) for its approximate metric differential at \(z\), a seminorm on \(T_zS^2\) describing its approximate first-order metric behavior. Under the usual qualitative convention, \(u\) is infinitesimally quasiconformal if there exists a finite constant \(Q_u\geq1\) such that, for almost every \(z\in S^2\),
\[
\sup_{\substack{v\in T_zS^2\\ |v|=1}}
\operatorname{ap\,md}u_z(v)
\leq
Q_u\inf_{\substack{v\in T_zS^2\\ |v|=1}}
\operatorname{ap\,md}u_z(v).
\]
The tangent norm is the round norm. This inequality permits the zero seminorm and excludes a nonzero seminorm with a nontrivial kernel. The constant is uniform in \(z\); the supplied statement specifies no numerical bound or uniformity across different maps or targets.

The question is: for every compact quasiconvex metric space \(X\) satisfying the stated local quadratic isoperimetric inequality and small-diameter null-homotopy hypothesis, is every harmonic map \(u:S^2\to X\) infinitesimally quasiconformal?

The meaning of harmonic is an unresolved qualification: the record does not specify the metric Sobolev energy or the variational condition defining harmonicity. Consequently, the question must retain the source's intended notion of harmonic map pending inspection of its definitions. It must not be restricted to maps that minimize energy on the entire sphere; the supplied assessment expressly distinguishes those maps from the general harmonic spheres asked about here.
\subsection{Short English statement}
Under the stated compactness, quasiconvexity, filling, and small-sphere homotopy assumptions on \(X\), must every harmonic sphere have bounded infinitesimal distortion almost everywhere? The intended energy and definition of harmonicity require source clarification. This is a statement-only draft requiring mathematical review.
\subsection{Sources}
\begin{itemize}
\item Supplied local UnsolvedMath registry record, permanent ID \(30006419\), title ``Infinitesimal Quasiconformality of Harmonic Spheres,'' source problem code \texttt{catalogue entry}, in the Oberwolfach Reports Open Problems collection. No individual proposer is identified in the record.
\item Supplied report citation: \emph{Differentialgeometrie im Großen (2026)}, Oberwolfach Reports, DOI \texttt{10.4171/owr/2025/30}, pp.\ 1591--1646. \url{https://doi.org/10.4171/owr/2025/30}. The supplied original excerpt asks whether \(u\) is infinitesimally quasiconformal but does not itself define \(u\).
\item Supplied literature source, labeled ``Meier--Vikman--Wenger Oberwolfach report'': \url{https://ems.press/content/serial-article-files/52249}.
\item Supplied literature source, \emph{Energy minimizing harmonic spheres}: \url{https://arxiv.org/abs/2503.08553}.
\item Literature assessment dated exactly \texttt{2026-08-22}: the supplied curation labels the problem open and states that the primary work proves infinitesimal quasiconformality for energy minimizers on the sphere while leaving the general harmonic-sphere question open. This is supplied dated curation, not an independently verified literature assessment. The registry's \texttt{corrected\_verified} statement label is likewise supplied metadata, not verification performed for this draft.
\end{itemize}
\end{document}
